\documentclass[11pt]{article}

\usepackage[a4paper,margin=1in]{geometry}
\usepackage[T1]{fontenc}
\usepackage[utf8]{inputenc}
\usepackage{lmodern}
\usepackage{microtype}
\usepackage{amsmath,amssymb,amsthm,mathtools}
\usepackage[colorlinks=true,linkcolor=blue,urlcolor=blue]{hyperref}
\usepackage[nameinlink,noabbrev]{cleveref}

\newtheorem{theorem}{Theorem}
\newtheorem{lemma}[theorem]{Lemma}
\newtheorem{proposition}[theorem]{Proposition}
\newtheorem{remark}[theorem]{Remark}

\newcommand{\Eerr}{\mathcal E}
\newcommand{\sens}{\mathrm{sens}}

\title{Low-Bandwidth \(\gamma\)-BIFR: Constant-Factor \(p\)-Flatness of the Collapsed Factorwise Proof}
\author{}
\date{}

\begin{document}
\maketitle

\section{Statement}

This note uses the notation from the BIFR paper.
The factorization is
\[
    E = B_\gamma^{(p)}C_\gamma^{(p)}
\]
and the native multi-participation error is
\[
    \Eerr(B_\gamma^{(p)},C_\gamma^{(p)})
    =
    \frac{\|B_\gamma^{(p)}\|_F}{\sqrt n}
    \sens_{k,b}(C_\gamma^{(p)}).
\]

The low-bandwidth policy in the current proof is
\[
    \eta := 1-\gamma = \frac{p}{\sqrt n},
    \qquad
    \gamma = 1-\frac{p}{\sqrt n},
\]
with
\[
    2\le p\le \frac{\sqrt n}{4}.
\]

The point of this note is narrow.
It proves that the collapsed factorwise low-bandwidth proof is \(p\)-flat: after multiplying the paper's RMSE and sensitivity estimates and simplifying to the asymptotic branch, the only remaining \(p\)-dependence is a bounded universal constant factor.
Thus this proof route cannot explain a polynomial, logarithmic, or asymptotic decrease in \(p\).

\section{Imported Bounds}

The current proof bounds the two factors separately.
We use the following two paper-level estimates.

\begin{proposition}[Normalized RMSE bound]
Under the low-bandwidth policy,
\[
    \frac{\|B_\gamma^{(p)}\|_F}{\sqrt n}
    \le
    C_R\sqrt{1+p^{2\eta}}
\]
for a universal constant \(C_R\).
\end{proposition}

\begin{proposition}[Strategy sensitivity bound]
Under the low-bandwidth policy,
\[
    \sens_{k,b}(C_\gamma^{(p)})
    \le
    C_S
    \left(
        \sqrt{k}\sqrt{1+\sqrt n\,p^{-2\eta}}
        +
        \sqrt{\frac{k}{b}}
        \left(1+\sqrt n\,p^{-\eta}\right)
    \right)
\]
for a universal constant \(C_S\).
\end{proposition}

The RMSE bound worsens with \(p\) through \(p^\eta\).
The sensitivity bound improves with \(p\) through \(p^{-\eta}\).
The flatness comes from multiplying these two effects.

\section{The Algebraic Cancellation}

Set
\[
    x := p^\eta = p^{p/\sqrt n}.
\]
Since \(p\ge2\), we have \(x\ge1\).
Since \(p\le\sqrt n/4\), we have \(\eta\le1/4\), so
\[
    x^2=p^{2\eta}\le p\le \frac{\sqrt n}{4}\le\sqrt n.
\]
Here the first inequality follows from \(2\eta\le1/2\le1\) and \(p\ge1\).

Multiplying the two imported estimates gives
\[
\begin{aligned}
    \Eerr(B_\gamma^{(p)},C_\gamma^{(p)})
    &\le
    C_R C_S
    \sqrt{1+x^2}
    \left(
        \sqrt{k}\sqrt{1+\sqrt n\,x^{-2}}
        +
        \sqrt{\frac{k}{b}}
        (1+\sqrt n\,x^{-1})
    \right).
\end{aligned}
\]
There are two product terms.
It is useful to name the collapsed factor envelope:
\[
    U_{n,k,b}(p)
    :=
    \sqrt{1+x^2}
    \left(
        \sqrt{k}\sqrt{1+\sqrt n\,x^{-2}}
        +
        \sqrt{\frac{k}{b}}
        (1+\sqrt n\,x^{-1})
    \right).
\]

For the first one,
\[
\begin{aligned}
    (1+x^2)(1+\sqrt n\,x^{-2})
    &=
    1+x^2+\sqrt n\,x^{-2}+\sqrt n  \\
    &\le
    2\sqrt n\,x^{-2}+2\sqrt n       \\
    &=
    2\sqrt n(1+x^{-2}),
\end{aligned}
\]
because \(x^2\le\sqrt n\) and \(1\le \sqrt n\,x^{-2}\).
Taking square roots,
\[
    \sqrt{1+x^2}\sqrt{1+\sqrt n\,x^{-2}}
    \le
    \sqrt2\, n^{1/4}\sqrt{1+x^{-2}}.
\]

For the second one, since \(x\ge1\),
\[
    \sqrt{1+x^2}\le \sqrt2\,x.
\]
Therefore
\[
    \sqrt{1+x^2}(1+\sqrt n\,x^{-1})
    \le
    \sqrt2(x+\sqrt n).
\]
The low-bandwidth assumptions imply \(n\ge64\) and \(x\le n^{1/4}\le\sqrt n\), so
\[
    \sqrt2(x+\sqrt n)
    \le
    2\sqrt2\,\sqrt n
    \le
    2\sqrt2\,\sqrt n\,\sqrt{1+x^{-2}}.
\]

\begin{lemma}[Two-sided flatness of the collapsed factor envelope]
Assume
\[
    k\ge1,
    \qquad
    b>0,
    \qquad
    2\le p\le\frac{\sqrt n}{4}.
\]
Then
\[
    \sqrt{k}\,n^{1/4}
    +
    \sqrt{\frac{kn}{b}}
    \le
    U_{n,k,b}(p)
    \le
    2\sqrt2
    \left(
        \sqrt{k}\,n^{1/4}
        +
        \sqrt{\frac{kn}{b}}
    \right).
\]
\end{lemma}

\begin{proof}
For the upper bound, the first product estimate and \(\sqrt{1+x^{-2}}\le\sqrt2\) give
\[
    \sqrt{1+x^2}\sqrt{1+\sqrt n\,x^{-2}}
    \le
    2n^{1/4}.
\]
The sharper intermediate estimate for the second product gives
\[
    \sqrt{1+x^2}(1+\sqrt n\,x^{-1})
    \le
    2\sqrt2\,\sqrt n.
\]
Since \(2\le2\sqrt2\), the desired upper comparison follows.
For the lower bound, the first product satisfies
\[
    \sqrt{1+x^2}\sqrt{1+\sqrt n\,x^{-2}}
    \ge
    \sqrt{\sqrt n}
    =
    n^{1/4},
\]
because the square of the left-hand side contains the term \(\sqrt n\).
The second product satisfies
\[
    \sqrt{1+x^2}(1+\sqrt n\,x^{-1})
    \ge
    x\cdot\frac{\sqrt n}{x}
    =
    \sqrt n.
\]
Multiplying these two lower bounds by \(\sqrt k\) and \(\sqrt{k/b}\), respectively, proves the lower comparison.
\end{proof}

Substituting the two upper product inequalities into the product bound gives
\[
    \Eerr(B_\gamma^{(p)},C_\gamma^{(p)})
    \le
    C
    \sqrt{1+x^{-2}}
    \left(
        \sqrt{k}\,n^{1/4}
        +
        \sqrt{\frac{kn}{b}}
    \right)
\]
for a universal constant \(C\).
Since \(x^{-2}=p^{-2p/\sqrt n}\), this is
\[
    \Eerr(B_\gamma^{(p)},C_\gamma^{(p)})
    \le
    C
    \sqrt{1+p^{-2p/\sqrt n}}
    \left(
        \sqrt{k}\,n^{1/4}
        +
        \sqrt{\frac{kn}{b}}
    \right).
\]

\section{Flatness Consequence}

\begin{theorem}[The collapsed factorwise low-bandwidth proof is \(p\)-flat]
Assume
\[
    k\ge1,
    \qquad
    b>0,
    \qquad
    2\le p\le \frac{\sqrt n}{4},
    \qquad
    \gamma=1-\frac{p}{\sqrt n}.
\]
The collapsed factor envelope satisfies
\[
    U_{n,k,b}(p)
    \asymp
    \sqrt{k}\,n^{1/4}
    +
    \sqrt{\frac{kn}{b}},
\]
with universal comparison constants.
Consequently, the collapsed factorwise proof route gives
\[
    \Eerr(B_\gamma^{(p)},C_\gamma^{(p)})
    \le
    C
    \sqrt{1+p^{-2p/\sqrt n}}
    \left(
        \sqrt{k}\,n^{1/4}
        +
        \sqrt{\frac{kn}{b}}
    \right),
\]
where \(C\) is universal.
Moreover
\[
    1
    \le
    \sqrt{1+p^{-2p/\sqrt n}}
    \le
    \sqrt2.
\]
Consequently, up to universal constants, this proof route is exactly the \(p\)-independent bound
\[
    \Eerr(B_\gamma^{(p)},C_\gamma^{(p)})
    =
    O\!\left(
        \sqrt{k}\,n^{1/4}
        +
        \sqrt{\frac{kn}{b}}
    \right).
\]
\end{theorem}

\begin{proof}
The two-sided comparison for \(U_{n,k,b}(p)\) is exactly the previous lemma.
The displayed p-sensitive estimate was proved in the previous section.
The residual factor satisfies
\[
    0 < p^{-2p/\sqrt n}\le1,
\]
so
\[
    1
    \le
    \sqrt{1+p^{-2p/\sqrt n}}
    \le
    \sqrt2.
\]
Thus the p-sensitive display and the p-independent display differ only by a universal constant.
In Big-O notation, the current factorwise proof has no asymptotic \(p\)-dependence.
\end{proof}

\begin{remark}[Boundary]
This proves flatness of the collapsed asymptotic upper-bound route.
It does not assert that the exact finite objective is flat in \(p\), and it does not assert that the uncollapsed finite-\(p\) envelope is flat.
\end{remark}

\end{document}
